\begin{tikzpicture}[
    x=1.25cm,
    y=0.95cm,
    wire/.style={
        draw=diagramfg,
        line cap=round,
        line width=0.8pt
    },
    crossing/.style={
        preaction={
            draw=diagrambg,
            line width=3.5pt,
            line cap=butt
        },
        draw=diagramfg,
        line width=0.8pt,
        line cap=round
    },
    bs/.style={
        draw=none,
        fill=diagramBS,
        rectangle,
        minimum width=5mm,
        minimum height=1.7mm,
        inner sep=0pt,
        line width=0.8pt
    },
    phase/.style={
        draw=diagramfg,
        text=diagramfg,
        fill=diagrambg,
        rectangle,
        rounded corners=1pt,
        minimum width=6mm,
        minimum height=6mm,
        line width=0.8pt,
    },
    meas/.style={
        draw=diagramfg,
        text=diagramfg,
        circle,
        rounded corners=1pt,
        minimum width=7mm,
        minimum height=7mm,
        inner sep=0pt,
        line width=0.8pt,
        fill=diagrambg,
        path picture={
            % semicircular dial
            \draw[draw=diagramfg,line width=0.8pt]
                ([xshift=-2.3mm,yshift=-1.1mm]
                    path picture bounding box.center)
                arc[start angle=180,end angle=0,radius=2.3mm];
            % needle with visible arrowhead
            \draw[draw=diagramfg,line width=0.9pt,
                -{Latex[length=1.5mm,width=1.2mm]}
            ]
                ([xshift=0mm,yshift=-1.0mm]
                    path picture bounding box.center)
                --
                ([xshift=1.7mm,yshift=2.6mm]
                    path picture bounding box.center);
        }
    },
    txt/.style={
        font=\small,
        text=diagramfg
    },
    epr/.style={
        draw=diagramaccent,
        dashed,
        line width=0.8pt
    },
    resource label/.style=epr bs/.style={
        draw=diagramfg,
        fill=diagrambg,
        rectangle,
        minimum width=6mm,
        minimum height=6mm,
        inner sep=0pt,
        line width=0.8pt
    },
    classical/.style={
        draw=diagramfg,
        double=diagrambg,
        double distance=1.1pt,
        line width=0.55pt,
        line cap=round,
        line join=round
    },
    classical arrow/.style={
        classical,
        -{Latex[length=1.8mm,width=1.4mm]},
        shorten >=1pt
    },
    wire arrow/.style={
        wire,
        -{Latex[length=1.8mm,width=1.4mm]},
        shorten >=1pt
    },
    disp/.style={
        draw=diagramfg,
        fill=diagrambg,
        rectangle,
        minimum size=6mm,
        inner sep=0pt,
        line width=0.8pt,
        text=diagramfg,
        rounded corners=1pt,
    }
]

\pgfdeclarelayer{foreground}
\pgfsetlayers{background,main,foreground}

% =================================================
% Rail order: (a,c,c',b,d,d')
% =================================================

\def\ya{1}
\def\yc{0}
\def\ycp{-1}

% =================================================
% Layout parameters
% Change these values to move gates without editing paths.
% =================================================

% Main horizontal stages
\def\xLabel{0}
\def\xSqueezing{1}
\def\xPhase{4.5}
\def\xMeasure{5.70}
\def\xResultLabel{6.5}
\def\xOutputEnd{5.2}
\def\xOutputLabel{5.25}

% EPR-resource geometry
% The "turn before" value is where the straight rail begins bending.
% The "port" value is where the bent segment reaches the BS edge.
% The "turn after" value is where the output reaches the original rail again.
\def\xEPRTurnBefore{1.85}
\def\xEPRPortBefore{2}
\def\xEPRBS{2.20}
\def\xEPRPortAfter{2.40}
\def\xEPRTurnAfter{2.55}
\def\yEPRUpper{-0.50}

% First beamsplitter layer geometry
\def\xFirstTurnBefore{2.90}
\def\xFirstPortBefore{3.05}
\def\xFirstBS{3.25}
\def\xFirstPortAfter{3.45}
\def\xFirstTurnAfter{3.60}
\def\yFirstUpper{0.50}


% Vertical port offset from each beamsplitter centre
\def\bsPortOffset{0.08}

% Classical feed-forward routing
\def\ffExitShort{0.25}
\def\ffExitLong{0.40}
\def\ffBusHeight{0.40}

%Vertical lines
\def\xGateLineOffset{0.6}
\def\xZeroLine{\xSqueezing-\xGateLineOffset}
\def\xOneLine{\xSqueezing+\xGateLineOffset}
\def\xTwoLine{{(\xEPRBS+\xFirstBS)/(2)}}
\def\xThreeLine{{(\xFirstBS+\xGateLineOffset)}}
\def\yLineEndUpper{1.2}
\def\yLineEndLower{-1.5}

% =================================================
% Left labels
% =================================================

\node[txt,anchor=east] at (\xLabel,\ya) {$A: \ket{\psi_{\mathrm{in}}}$};
\node[txt,anchor=east] (cLabel) at (\xLabel,\yc) {$B: \ket{0}$};
\node[txt,anchor=east] (cpLabel) at (\xLabel,\ycp) {$C: \ket{0}$};

% =================================================
% Vertical step lines
% =================================================
\draw[dotted, color=diagramfg] (\xZeroLine,\yLineEndUpper) -- (\xZeroLine,\yLineEndLower);
\draw[dotted, color=diagramfg] (\xOneLine,\yLineEndUpper) -- (\xOneLine,\yLineEndLower);
\draw[dotted, color=diagramfg] (\xTwoLine,\yLineEndUpper) -- (\xTwoLine,\yLineEndLower);
\draw[dotted, color=diagramfg] (\xThreeLine,\yLineEndUpper) -- (\xThreeLine,\yLineEndLower);


\node[txt,anchor=south] at (\xZeroLine,\yLineEndUpper) {$0$};
\node[txt,anchor=south] at (\xOneLine,\yLineEndUpper) {$1$};
\node[txt,anchor=south] at (\xTwoLine,\yLineEndUpper) {$2$};
\node[txt,anchor=south] at (\xThreeLine,\yLineEndUpper) {$3$};
% =================================================
% Squeezing
% =================================================




\node[phase] (Sp) at (\xSqueezing,\yc) {$S(r,0)$};
\node[phase] (Sx) at (\xSqueezing,\ycp) {$S(r,\pi)$};

\draw[wire] (0.12,\yc)  -- (Sp.west);
\draw[wire] (0.12,\ycp) -- (Sx.west);


% =================================================
% EPR resources
% =================================================
% ---------- upper EPR beamsplitter between c and c' ----------

% left straight parts
\draw[wire] (Sp.east) -- (\xEPRTurnBefore,\yc);
\draw[wire] (Sx.east) -- (\xEPRTurnBefore,\ycp);

% gaps
%\draw[crossing] (3.50,\ycp) -- (3.60,\ycp);
%\draw[wire] (0.12,\ycp) -- (\xEPRTurnBefore,\ycp);


% bend into the BS
\draw[wire] (\xEPRTurnBefore,\yc)  -- (\xEPRPortBefore,{\yEPRUpper+\bsPortOffset});
\draw[wire] (\xEPRTurnBefore,\ycp) -- (\xEPRPortBefore,{\yEPRUpper-\bsPortOffset});

% bend back out to the original rails
\draw[wire] (\xEPRPortAfter,{\yEPRUpper+\bsPortOffset}) -- (\xEPRTurnAfter,\yc);
\draw[wire] (\xEPRPortAfter,{\yEPRUpper-\bsPortOffset}) -- (\xEPRTurnAfter,\ycp);



% beamsplitter
\node[bs] (EPRup) at (\xEPRBS,\yEPRUpper) {};

% framed label

\node[
    draw=diagramaccent!75!black,
    rounded corners=2pt,
    line width=0.9pt,
    inner xsep=8pt,
    inner ysep=4pt,
    fit=(cLabel)(cpLabel)(EPRup),
    label={[text=diagramaccent!75!black,font=\small]above:EPR}
    ] (upperEPRbox) {};
% ---------- lower EPR beamsplitter between d and d' ----------



% =================================================
% First layer: B_ac and B_bd
% =================================================

% Inputs to B_ac
\draw[wire] (0.12,\ya) -- (\xFirstTurnBefore,\ya) -- (\xFirstPortBefore,{\yFirstUpper+\bsPortOffset});
\draw[wire] (\xEPRTurnAfter,\yc) -- (\xFirstTurnBefore,\yc) -- (\xFirstPortBefore,{\yFirstUpper-\bsPortOffset});


% Outputs of B_ac: return to rails a,c
\coordinate (acup) at (\xFirstTurnAfter,\ya);
\coordinate (aclo) at (\xFirstTurnAfter,\yc);



\node[bs] (Bac) at (\xFirstBS,\yFirstUpper) {};

%\node[txt,above=2mm of Bac] {$B_{01}$};



% =================================================
% Phase rotations
% =================================================

\node[phase] (R3) at (\xPhase,\ya) {$R(\theta_3)$};
\node[phase] (R1) at (\xPhase,\yc) {$R(\theta_1)$};

\draw[wire] (\xFirstPortAfter,{\yFirstUpper+\bsPortOffset}) -- (acup) -- (R3.west);
\draw[wire] (\xFirstPortAfter,{\yFirstUpper-\bsPortOffset}) -- (aclo) -- (R1.west);
% =================================================
% Measurements
% =================================================

\node[meas] (M3) at (\xMeasure,\ya) {};
\node[meas] (M1) at (\xMeasure,\yc) {};

\draw[wire] (R3.east) -- (M3.west);
\draw[wire] (R1.east) -- (M1.west);


\node[txt,anchor=west] at (\xResultLabel,\ya) {$m_3$};
\node[txt,anchor=west] at (\xResultLabel,\yc) {$m_1$};

% =================================================
% Surviving outputs
% =================================================
% Displacement gate on c'



% continue on the rails
\draw[wire arrow] (\xEPRTurnAfter,\ycp) -- (\xOutputEnd,\ycp);


% Dedicated routing positions
\coordinate (M3exit) at ($(M3.east)+(\ffExitLong,0)$);
\coordinate (M1exit) at ($(M1.east)+(\ffExitLong,0)$);


% Measurement outcomes
\draw[classical] (M3.east) -- (M3exit) -- ($(\xResultLabel, \ya)-(0.25,0)$);

\draw[classical] (M1.east) -- (M1exit)  -- ($(\xResultLabel, \yc)-(0.25,0)$);



\node[
    txt,
    anchor=west
] at (\xOutputLabel,\ycp)
    {$V_{0}(\theta_1,\theta_3)\ket{\psi_{\mathrm{in}}}$};


% =================================================
% Bottom caption
% =================================================

\node[txt,align=center] at (2.85,-1.75)
{
    one-mode gate teleportation
};

\end{tikzpicture}
